\section*{الفصل \ref{ch:g11:dna-replication} --- تضاعف DNA}

\begin{solution}{exo:g11:dna-replication:1}
يعمل كل شريط أصلًا يبنى عليه شريط متمم جديد. ولا يمكن ذلك إلا لأن A
تزاوج T وG تزاوج C: فمتتالية أحد الشريطين تثبّت متتالية الآخر، ومن ثم
يحمل الشريط وحده المعلومة كلها.
\end{solution}

\begin{solution}{exo:g11:dna-replication:2}
النصف المحافظ: يحتفظ كل جزيء بنت بشريط أبوي ويكسب شريطًا جديدًا. أما
المحافظ: فيبقى الجزيء الأب كاملًا ويبنى بجانبه شريط مزدوج جديد كله.
\end{solution}

\begin{solution}{exo:g11:dna-replication:3}
المرحلة S من الدورة الخلوية؛ وتتضاعف كمية \omterm{def:g10:universal-dna:information}{DNA} في \omterm{def:g10:cells-common-unit:organelle}{النواة}.
\end{solution}

\begin{solution}{exo:g11:dna-replication:4}
\omterm{prop:g11:dna-replication:forks}{بوليميراز DNA} يضيف إلى شريط نامٍ \omterm{def:g10:universal-dna:nucleotide}{النوكليوتيدات} المتممة للأصل، واحدًا
واحدًا. أما \omterm{prop:g11:dna-replication:forks}{شوكة التضاعف} فهي النقطة المتحركة التي يفصل عندها الشريطان
الأبويان ويبنى الشريطان الجديدان.
\end{solution}

\begin{solution}{exo:g11:dna-replication:5}
\texttt{CTAATGTCCG}، قاعدة تحت قاعدة.
\end{solution}

\begin{solution}{exo:g11:dna-replication:6}
المحافظ، الجيل 1: نطاق ثقيل ونطاق خفيف، ولا هجين --- لكن ما رئي إلا
نطاق هجين. والمبدد، الجيل 2: نطاق واحد على ربع المسافة من الخفيف إلى
الثقيل --- لكن ما رئي هو نطاقان متمايزان. فكل نموذج يتنبأ بنطاقات لم
تشاهد.
\end{solution}

\begin{solution}{exo:g11:dna-replication:7}
$2^5 = 32$ جزيئًا، 2 هجينان و30 خفيفًا: أي نطاق خفيف خمسة عشر ضعف
الهجين، ولا نطاق ثقيل.
\end{solution}

\begin{solution}{exo:g11:dna-replication:8}
الجيل 0: نطاق خفيف واحد. والجيل 1: نطاق هجين واحد. والجيل 2: نطاقان
هجين وثقيل متساويا الشدة --- أي صورة معكوسة للتجربة الأصلية.
\end{solution}

\begin{solution}{exo:g11:dna-replication:9}
$\num{2.5e8}/50 = \num{5e6}\,\unit{s}$، أي نحو 58 يومًا --- وهو أطول
بكثير من 8 ساعات. فوجب نسخ \omterm{def:g10:universal-dna:information}{الصبغي} من منشآت كثيرة معًا، يفتح كل منها
شوكتين.
\end{solution}

\begin{solution}{exo:g11:dna-replication:10}
شريطا $\num{3.2e9}$ زوج قواعد، مرتين: أي $\num{6.4e9}$ زوج قواعد،
ومنه $\num{6.4e9}$ \omterm{def:g10:universal-dna:nucleotide}{نوكليوتيد} جديد وصل بعضه ببعض (واحد لكل زوج قواعد
منسوخ). وعند $10^{-9}$ لكل \omterm{def:g10:universal-dna:nucleotide}{نوكليوتيد}، نحو 6 أخطاء.
\end{solution}

\begin{solution}{exo:g11:dna-replication:11}
فتح فقاعة يكشف أصلًا على جانبي المنشأ؛ وكل شوكة تنسخ مبتعدةً عن
المنشأ، فتسير الشوكتان في اتجاهين متعاكسين وتغطيان معًا المنطقة كلها.
وحين تلتقي فقاعتان تندمج شوكتاهما وتتصل الشرائط البنات طرفًا بطرف في
جزيئات متصلة.
\end{solution}

\begin{solution}{exo:g11:dna-replication:12}
تكف البكتيريا عن الانقسام (فلا تضاعف ولا انقسام). ويكف الجرح عن
الالتئام، لأن إصلاحه يحتاج إلى انقسامات خلوية. أما \omterm{def:g10:cells-common-unit:cell}{الخلية} العصبية فلا
تتأثر: فهي لم تعد \omterm{def:g11:dna-replication:replication}{تضاعف DNA} فيها. ودواء يوقف \omterm{def:g10:cells-common-unit:cell}{الخلايا} السريعة الانقسام
ويعفي غير المنقسمة يوحي بعلاج ضد البكتيريا أو ضد \omterm{def:g10:cells-common-unit:cell}{خلايا} سرطانية
منقسمة.
\end{solution}

\begin{solution}{exo:g11:dna-replication:13}
الشريطان الأصليان لا يهدمان قط: ففي كل جيل يزاوج كل منهما من جديد
شريطًا خفيفًا جديدًا، فيوجد دائمًا جزيئان هجينان بالضبط. وبعد 10
أجيال: $2/2^{10} = 2/1024
\approx 0.2\%$ --- موجود لكنه أخفت من أن يرى.
\end{solution}

\begin{solution}{exo:g11:dna-replication:14}
$\num{6.4e9} \times 10^{-5} = \num{64000}$ خطأ في كل انقسام. ومع عشرين
ألف \omterm{def:g10:universal-dna:gene}{مورّثة}، كانت مئات \omterm{def:g10:universal-dna:gene}{المورّثات} لتتضرر في كل انقسام؛ وبعد بضعة
انقسامات كانت \omterm{def:g10:cells-common-unit:cell}{الخلايا} لتحمل عيوبًا مميتة وينقرض النسل.
\end{solution}

\begin{solution}{exo:g11:dna-replication:15}
تبدأ جولة تضاعف جديدة عند المنشأ قبل أن تنتهي الجولة السابقة: فيحمل
\omterm{def:g10:universal-dna:information}{الصبغي} عدة أزواج متداخلة من الشوكات في آن واحد. فيتلقى كل انقسام
عندئذ صبغيًا منسوخًا جزئيًا من قبل للانقسام التالي، وتتلو الانقسامات
بعضها بعضًا أسرع مما يستغرقه نسخ واحد.
\end{solution}

\begin{solution}{pb:g11:dna-replication:1}
\textbf{1.} النيتروجين 16\% من الكتلة وهو أثقل 7\%:
$0.16 \times 0.07 \approx 1.1\%$. وطريقة التدرج تفصل بين كثافتين
تختلفان بكسر من المئة، فنسبة 1\% تعطي نطاقين متباعدين تباعدًا واضحًا.

\textbf{2.} حتى يكون \omterm{def:g10:universal-dna:information}{DNA} كله ثقيلًا في الجوهر، لا نصفه أو ثلاثة
أرباعه: فبعد 14 جيلًا يبقى من الشرائط الخفيفة الأصلية أقل من جزء في
$10^4$.

\textbf{3.} في منتصف المسافة بين النطاقين بالضبط: فكثافته هي الوسطي.

\textbf{4.} البكتيريا الواحدة تحمل بضعة فمتوغرامات من \omterm{def:g10:universal-dna:information}{DNA}؛ والصورة
تحتاج إلى ميكروغرامات، ومنه إلى مليارات \omterm{def:g10:cells-common-unit:cell}{الخلايا}.

\textbf{5.} بأن يبقيا مزرعة في وسط ثقيل شاهدًا ويتحققا من بقاء نطاقها
ثقيلًا، وبأن يتحققا من النطاق الخفيف لبكتيريا نمت في وسط خفيف طوال
الوقت.

\textbf{6.} النصف المحافظ: الجيل 1، نطاق هجين واحد؛ والجيل 2، نطاقان
هجين وخفيف متساويان.

\textbf{7.} المحافظ: الجيل 1، نطاقان ثقيل وخفيف متساويان؛ والجيل 2،
ثقيل واحد مقابل ثلاثة خفيفة.

\textbf{8.} المبدد: الجيل 1، نطاق واحد في الموضع الهجين؛ والجيل 2،
نطاق واحد على ربع المسافة من الخفيف إلى الثقيل.

\textbf{9.} يستبعد النموذج المحافظ، وقد تنبأ بألا نطاق هجين. أما
النصف المحافظ والمبدد فيتنبآن كلاهما بالنطاق الهجين الواحد ويبقيان.

\textbf{10.} ينجو النصف المحافظ: فهو وحده يتنبأ بنطاقين منفصلين. أما
النموذج المبدد فيتنبأ بنطاق واحد ينزلق نحو الخفة، ولا ينقسم قط.

\textbf{11.} $2^n$ جزيء، منها 2 بالضبط تحمل شريطًا أصليًا.

\textbf{12.} الخفيف إلى الهجين $= (2^n - 2) : 2$: فالجيل 3، $3:1$؛
والجيل 4، $7:1$؛ والجيل 6، $31:1$.

\textbf{13.} النسبة الهجينة $2/2^n < 0.05$ حين $2^n > 40$: أي من
الجيل 6 ($2/64 \approx 3\%$).

\textbf{14.} نطاقان، ثقيل وخفيف، بمقدارين متساويين: فالجزيء الهجين
شريط ثقيل زائد شريط خفيف. وهذا يختبر النموذج المبدد، وهو يتنبأ بأن كل
شريط مفرد ذو كثافة وسطى، فيستبعده استبعادًا مستقلًا.

\textbf{15.} انطلاقًا من الخفيف، يعطي جيل واحد في وسط ثقيل نطاقًا
هجينًا واحدًا، ويعطي جيلان هجينًا وثقيلًا: أي النمط نفسه معكوسًا، بما
يبين أن النتيجة لا تتوقف على أي النظيرين هو «القديم».

\textbf{16.} $\num{4.6e6}/(2 \times 1000) = \qty{2300}{s}$، أي نحو
38 دقيقة.

\textbf{17.} 77 دقيقة. وعلى حلقة يكون الشريطان مكشوفين على جانبي أي
فتحة؛ وشوكتان تغادران منشأً واحدًا في اتجاهين متعاكسين هما أبسط طريقة
لتغطية الحلقة كلها.

\textbf{18.} المنشأ الواحد ينسخ $2 \times 50 \times 8 \times 3600 =
\num{2.88e6}$ زوج قواعد في 8 ساعات؛ و$\num{6.4e9}/\num{2.88e6}
\approx 2200$ منشأ على أقل تقدير (وفي الواقع عشرات الألوف، تعمل في
أوقات مختلفة).

\textbf{19.} البكتيريا: $\num{9.2e6} \times 10^{-9} \approx 0.01$ خطأ
في كل تضاعف؛ \omterm{def:g10:cells-common-unit:cell}{والخلية} البشرية: نحو 6. ومعظم المجين البشري ليس \omterm{def:g10:universal-dna:gene}{مورّثات}،
فمعظم الأخطاء تقع حيث لا تغير شيئًا، وللخلية نسختان من كل \omterm{def:g10:universal-dna:gene}{مورّثة}.

\textbf{20.} النطاقان عند الجيل الثاني --- هجين وخفيف بمقدارين
متساويين --- أثبتا أن التضاعف نصف محافظ؛ والشوكتان تنسخان \omterm{def:g10:universal-dna:information}{الصبغي}
الجرثومي في نحو 38 دقيقة، أي زمن جيل البكتيريا تقريبًا، فالنسخ يضبط
وتيرة الانقسام.
\end{solution}
